% Energy surfaces of the MomC model as functions of momentum (cases 1 and 2).
% Input: MOM<case>_1.dat, MOM<case>_2.dat (lower/upper surface), created by
% momentum_surfaces.m.

\documentclass{standalone}

\usepackage{amsmath}
\usepackage{pgfplots}
\usetikzlibrary{matrix}
\usetikzlibrary{patterns}
\usepgfplotslibrary{fillbetween}

\begin{document}

% cases 1 and 2
\begin{tabular}{cc}
	%
	\begin{tikzpicture}
		%
		\begin{axis}[
			width = 16cm,
			height = 8cm,
			scale = .7,
			xmin = -6,
			xmax = 6,
			ymin = -1,
			ymax = 4,
			xlabel = {Momentum (eV/c)},
			ylabel = {Energy (meV)},
			grid = both,
			]
			\addplot[thick] table {MOM1_1.dat};
			\label{lambda2}
			\addplot[thick, red] table {MOM1_2.dat};
			\label{lambda1}
		\end{axis}
		%
		\matrix
			[
			matrix of nodes,
			anchor = south east,
			fill = white,
			draw,
			inner sep = 2pt,
			node font = \scriptsize,
			column 1/.style={nodes={align=left, text width=8mm}},
			column 2/.style={nodes={align=left, text width=4mm}},
			xshift = -1mm,
			yshift = 1mm]
		at (current axis.south east) {
			\ref{lambda1} & $\lambda_1$\\[1mm]
			\ref{lambda2} & $\lambda_2$\\[1mm]
			};
	\end{tikzpicture}
	%
	&
	%
	\begin{tikzpicture}
		%
		\begin{axis}[
			width = 16cm,
			height = 8cm,
			scale = .7,
			xmin = -6,
			xmax = 6,
			ymin = -1,
			ymax = 4,
			xlabel = {Momentum (eV/c)},
			ylabel = {Energy (meV)},
			grid = both,
			]
			\addplot[thick] table {MOM2_1.dat};
			\addplot[thick, red] table {MOM2_2.dat};
		\end{axis}
		%
	\end{tikzpicture}
	%
\end{tabular}

\end{document}
